% =====================================================================
%  Seance 1 -- A. page de garde, plan, ML fondamental (frames 1--13)
% =====================================================================
\section{Séance 1 --- Introduction au ML et aux infrastructures}
\chapterslide{Séance 1}{Introduction à l'apprentissage automatique et aux infrastructures}
{4 h --- 2 h cours / 2 h TP --- Partie 1 du syllabus EPD-GIT-5534}
{\textbf{Objectifs :} situer les trois familles de ML $\bullet$ nommer les composants d'une
infrastructure et leurs indicateurs de santé (SLI/SLO/SLA, MTBF/MTTR) $\bullet$
formuler un problème opérationnel en termes de ML $\bullet$ explorer un jeu de
métriques réelles (NOSYS) et produire une fiche « état des lieux » chiffrée.}

\begin{frame}{Plan de la séance (4 h)}
\begin{columns}[t]
\begin{column}{0.48\textwidth}
\begin{kbox}[Cours magistral --- 2 h]
\begin{enumerate}
\item \textbf{Notions fondamentales du ML} (50 min) : les trois familles, le vocabulaire.
\item \textbf{Anatomie d'une infrastructure} (40 min) : couches, métriques, SLI/SLO/SLA, MTBF/MTTR.
\item \textbf{Le croisement : AIOps} (20 min) : pourquoi le ML change la supervision.
\item \textbf{Lancement du dossier NOSYS} (10 min).
\end{enumerate}
\end{kbox}
\end{column}
\begin{column}{0.48\textwidth}
\begin{ebox}[Travaux pratiques --- 2 h]
\textbf{« Cartographier la plateforme »} sur les vraies métriques
\file{nosys\_metriques.csv} :
\begin{enumerate}
\item B1 --- chargement et vérifications (30 min)
\item B2 --- inventaire de qualité (20 min)
\item B3 --- visualiser un serveur sur 7 jours (30 min)
\item B4 --- statistiques par serveur (40 min)
\item B5 --- livrable : fiche « état des lieux » (20 min)
\end{enumerate}
\end{ebox}
\end{column}
\end{columns}
\vspace{1mm}
\begin{center}\small
Fil conducteur des 6 séances : \textbf{NOSYS}, exploitant d'un centre de données
et d'un réseau 4G/5G à Douala.
\end{center}
\end{frame}

\begin{frame}{Pourquoi le ML change le métier d'exploitant}
\footnotesize%
\begin{columns}[t]
\begin{column}{0.52\textwidth}
\begin{ebox}[Le problème NOSYS (données réelles du dossier)]
Les journaux du mois (\file{nosys\_logs.csv}) contiennent \textbf{100\,000 lignes}
d'authentification. Une intrusion brute force y est enfouie :
\textbf{125 échecs d'authentification concentrés} sur 2 heures de nuit,
sur les comptes \py{root}, \py{admin}, \py{inconnu-*}.
\begin{itemize}
\item Baseline normale : $\approx$ 30 échecs/jour sur tout le parc.
\item Personne ne lit 100\,000 lignes : \textbf{l'œil humain rate l'attaque}.
\end{itemize}
\end{ebox}
\begin{kbox}[La promesse de cette UE]
Un modèle qui lit les données en continu, apprend ce qui est \emph{normal},
et alerte \emph{avant} la panne --- pas 3 heures après.
\end{kbox}
\end{column}
\begin{column}{0.44\textwidth}
\begin{wbox}[Ce que le ML ne remplace pas]
\begin{itemize}
\item Le \textbf{bon sens opérationnel} : un modèle alerte, un humain décide.
\item Les \textbf{règles simples} quand elles suffisent (un disque plein, on le voit).
\item La \textbf{qualité des données} : garbage in, garbage out.
\end{itemize}
\end{wbox}
\vspace{2mm}
\begin{center}\small
Objectifs pédagogiques du syllabus couverts aujourd'hui :\\
\#1 concepts fondamentaux $\bullet$ \#2 usages dans les infrastructures.
\end{center}
\end{column}
\end{columns}
\end{frame}

\begin{frame}{Qu'est-ce que l'apprentissage automatique ?}
\begin{block}{Définition de travail}
\textbf{Apprendre} une fonction de prédiction ou de structure
$f$ à partir de \textbf{données}, plutôt que de l'écrire à la main :
\[
\text{données } (x, y) \;\longrightarrow\; \text{apprentissage}
\;\longrightarrow\; f \;\longrightarrow\; \text{prédictions sur du \textbf{nouveau}}
\]
\end{block}
\begin{itemize}
\item \textbf{Programmation classique} : l'humain écrit les règles
(\py{if cpu > 90: alerte()}).
\item \textbf{ML} : l'humain fournit des exemples, la machine \emph{apprend} les règles
--- y compris celles qu'on ne saurait pas écrire (quel pattern de latence annonce une panne ?).
\item Le critère d'apprentissage (fonction de coût) et la mesure de généralisation
sont les \textbf{deux piliers} de tout modèle : on les verra à chaque séance.
\end{itemize}
\begin{kbox}[À retenir dès maintenant]
Un modèle vaut ce que vaut sa performance \textbf{sur des données qu'il n'a jamais vues}
--- jamais sa performance sur les données d'entraînement.
\end{kbox}
\end{frame}

\begin{frame}{Les trois familles de l'apprentissage}
\footnotesize%
\begin{center}
\begin{tikzpicture}[node distance=4mm,
  box/.style={rectangle,rounded corners=2mm,draw,align=center,font=\scriptsize,minimum height=9mm}]
\node[box,draw=byBlue,fill=byBlue!8,text width=3.6cm] (sup) {\textbf{Supervisé}\\données étiquetées $(x,y)$\\régression \fcolorbox{byDeep}{byCyan!20}{\scriptsize ou} classification};
\node[box,draw=byCyan,fill=byCyan!8,right=of sup,text width=3.6cm] (nonsup) {\textbf{Non supervisé}\\pas d'étiquette, chercher une\\structure : clustering,\\détection d'anomalies};
\node[box,draw=byOrange,fill=byOrange!8,right=of nonsup,text width=3.6cm] (renf) {\textbf{Par renforcement}\\un agent agit, reçoit une\\récompense, ajuste\\sa stratégie};
\draw[-{Stealth},thick,byDeep] (sup) -- node[above,font=\tiny]{étiquettes} (nonsup);
\end{tikzpicture}
\end{center}
\vspace{1mm}
\begin{columns}[t]
\begin{column}{0.32\textwidth}
\begin{ebox}[En infrastructure]
Prédire la latence à $h{+}1$ ;\\
classer une session « attaque / légitime » ;\\
estimer la durée d'un incident.
\end{ebox}
\end{column}
\begin{column}{0.32\textwidth}
\begin{ebox}[En infrastructure]
Regrouper les serveurs par profil de charge ;\\
détecter les fenêtres « anormales » ;\\
segmenter les applications.
\end{ebox}
\end{column}
\begin{column}{0.32\textwidth}
\begin{ebox}[En infrastructure]
Piloter le refroidissement d'une salle ;\\
ajuster dynamiquement la\\limite de débit d'un client.
\end{ebox}
\end{column}
\end{columns}
\vspace{1mm}
\begin{center}\small
Cette UE couvre en profondeur les \textbf{deux premières familles} ;
le renforcement est un aperçu culturel.
\end{center}
\end{frame}

\begin{frame}{Supervisé (1/2) : la régression}
\footnotesize%
\begin{block}{Principe}
La cible $y$ est \textbf{continue}. On apprend $f(x) \approx y$, par exemple une forme linéaire
\[
\hat{y} = \beta_0 + \beta_1 x_1 + \dots + \beta_p x_p
\]
en minimisant l'erreur quadratique moyenne (MSE) :
\[
\mathrm{MSE} = \frac{1}{n}\sum_{i=1}^{n}\bigl(y_i - \hat{y}_i\bigr)^2
\]
\end{block}
\begin{columns}[t]
\begin{column}{0.48\textwidth}
\begin{ebox}[Exemples « infrastructure »]
\begin{itemize}
\item Prédire la \textbf{latence} de demain 14 h--15 h à partir de la charge.
\item Prédire la \textbf{consommation} (kWh) d'une salle la semaine prochaine
--- notre parc a consommé \textbf{11\,614 kWh} en 30 jours : le modèle du TP de la
séance 4 partira de ce chiffre.
\item Estimer la \textbf{durée restante} avant saturation du stockage.
\end{itemize}
\end{ebox}
\end{column}
\begin{column}{0.48\textwidth}
\begin{kbox}[Ce qu'on apprendra (séances 4 et 11)]
\begin{itemize}
\item Régression linéaire : interpréter un $\beta$ (« +1 \% de CPU $\Rightarrow$ +x ms »).
\item Séries temporelles (ARIMA) : les métriques d'infra sont \textbf{correlées dans le temps},
on ne peut pas les traiter comme un tableau aléatoire.
\end{itemize}
\end{kbox}
\end{column}
\end{columns}
\end{frame}

\begin{frame}{Supervisé (2/2) : la classification}
\footnotesize%
\begin{block}{Principe}
La cible $y$ est \textbf{discrète} (une classe). Le modèle sort une \textbf{probabilité}
$P(y = k \given x)$ et une règle de décision (seuil).
\[
P(\text{incident} \given \text{CPU, latence, heure}) = 0{,}87
\;\Rightarrow\; \text{alerte si } P > 0{,}5
\]
\end{block}
\begin{columns}[t]
\begin{column}{0.48\textwidth}
\begin{ebox}[Exemples « infrastructure »]
\begin{itemize}
\item \textbf{Incident dans les 30 min ?} (cible binaire, construite au TP2)
\item \textbf{Session attaque / légitime} (séance 5)
\item \textbf{Triage} d'un ticket : réseau / matériel / applicatif
\end{itemize}
\end{ebox}
\begin{wbox}[Le piège des classes rares]
Les incidents sont rares (16 blocs sur un mois) : la simple
\textbf{exactitude} est trompeuse. Un modèle qui dit toujours
« pas d'incident » a 99,98 \% d'exactitude et \textbf{zéro valeur}.
$\Rightarrow$ précision, rappel, F1 (séance 3).
\end{wbox}
\end{column}
\begin{column}{0.48\textwidth}
\begin{kbox}[Vocabulaire de décision]
\begin{itemize}
\item \textbf{Faux positif} : alerte injustifiée $\Rightarrow$ fatigue d'alerte du NOC.
\item \textbf{Faux négatif} : panne non détectée $\Rightarrow$ SLA cassé.
\item Les deux coûts sont \textbf{asymétriques} : c'est un choix métier, pas technique.
\end{itemize}
\end{kbox}
\end{column}
\end{columns}
\end{frame}

\begin{frame}{Non supervisé : clustering et détection d'anomalies}
\footnotesize%
\begin{columns}[t]
\begin{column}{0.48\textwidth}
\begin{block}{Clustering}
Regrouper les observations \textbf{sans étiquette}. Ex. : $k$-means
partitionne les serveurs selon (CPU, RAM, latence) en $k$ profils.
\begin{itemize}
\item Sert à \textbf{comprendre} avant de modéliser.
\item Séance 10 (cours de ML général) / séance 5 ici pour regrouper des attaques.
\end{itemize}
\end{block}
\begin{ebox}[Sur NOSYS]
Classer nos \textbf{8 serveurs} par profil de charge : avec les CPU moyens
(14,4 à 24,6 \%), les profils sont proches --- le clustering devient intéressant
quand on agrège par \emph{application}, pas par machine.
\end{ebox}
\end{column}
\begin{column}{0.48\textwidth}
\begin{block}{Détection d'anomalies}
Apprendre la distribution du \textbf{normal}, puis mesurer l'improbabilité :
\[
\text{score}(x) \text{ grand}
\;\Longleftrightarrow\;
x \text{ improbable sous } p_{\text{norm}}
\]
\begin{itemize}
\item C'est \textbf{le} problème central de la supervision.
\item Outils de la séance 3 : $z$-score, IQR, isolation forest, LOF.
\end{itemize}
\end{block}
\begin{kbox}[Pourquoi pas seulement des seuils ?]
Parce que « normal » \textbf{dépend de l'heure} : CPU 30 \% à 14 h est normal,
CPU 30 \% à 3 h du matin pour srv-01 est suspect. Un seuil fixe ne voit pas cette différence.
\end{kbox}
\end{column}
\end{columns}
\end{frame}

\begin{frame}{Par renforcement : l'aperçu culturel}
\footnotesize%
\begin{block}{Principe}
Un \textbf{agent} observe un état, \textbf{agit}, reçoit une \textbf{récompense},
et ajuste sa politique pour maximiser la récompense cumulée.
Pas de données étiquetées : les \emph{essais} créent l'expérience.
\end{block}
\begin{columns}[t]
\begin{column}{0.5\textwidth}
\begin{ebox}[L'exemple qui a tout déclenché]
Les datacenters de Google : un agent de renforcement pilotant la
climatisation a réduit l'énergie de refroidissement d'environ \textbf{40 \%}
(DeepMind, 2016). La récompense : PUE (efficacité énergétique) minimisée.
\end{ebox}
\begin{ebox}[Autres usages « infra »]
\begin{itemize}
\item Routage dynamique dans un réseau maillé.
\item Ordonnancement de tâches sur un parc hétérogène.
\end{itemize}
\end{ebox}
\end{column}
\begin{column}{0.46\textwidth}
\begin{wbox}[Pourquoi un aperçu seulement]
\begin{itemize}
\item Nécessite un \textbf{simulateur} fidèle de l'infrastructure
(on ne fait pas d'essais/erreurs en production !).
\item Hors du périmètre de cette UE : cité pour que vous
\textbf{sachiez le nommer} et reconnaître ses usages.
\end{itemize}
\end{wbox}
\begin{center}\small
En 2026, la majorité des usages ML \emph{matures} en infrastructure
sont supervisés ou non supervisés.
\end{center}
\end{column}
\end{columns}
\end{frame}

\begin{frame}{Exemple : classer les problèmes de NOSYS}
\begin{columns}[t]
\begin{column}{0.48\textwidth}
\begin{abox}[À vous de jouer --- 4 min]
Classer chaque problème dans la bonne famille (et sous-famille) :
\begin{enumerate}
\item Prévoir le trafic agrégé de la semaine prochaine.
\item Détecter les fenêtres de métriques « pas comme les autres ».
\item Prédire si un incident survient dans les 30 min.
\item Regrouper les applications par profil de consommation.
\item Choisir automatiquement quand lancer une tâche de sauvegarde.
\end{enumerate}
\end{abox}
\end{column}
\begin{column}{0.48\textwidth}
\begin{solbox}[Corrigé]
\begin{enumerate}
\item Supervisé --- \textbf{régression} (série temporelle : séance 4).
\item Non supervisé --- \textbf{anomalies} (séance 3).
\item Supervisé --- \textbf{classification} binaire (séances 3--4).
\item Non supervisé --- \textbf{clustering}.
\item Renforcement --- l'agent apprend une politique d'ordonnancement
(aperçu, hors UE).
\end{enumerate}
\textbf{Erreur fréquente} : le n°2 n'est pas supervisé --- il n'existe pas
d'étiquette « anormal » au moment de l'apprentissage.
\end{solbox}
\end{column}
\end{columns}
\end{frame}

\begin{frame}{Le vocabulaire qui reviendra à chaque séance}
\begin{columns}[t]
\begin{column}{0.49\textwidth}
\begin{kbox}[Les mots et leur sens exact]
\begin{description}
\item[Variable explicative ($x$)] ce qu'on observe : CPU, latence, heure\ldots
\item[Variable cible ($y$)] ce qu'on prédit : incident, charge, classe.
\item[Jeu d'entraînement] données qui \textbf{fabriquent} le modèle.
\item[Jeu de test] données \textbf{jamais vues} qui mesurent la généralisation.
\item[Surapprentissage] le modèle a mémorisé le bruit du train ---
excellent sur le passé, mauvais demain.
\item[Pipeline] la chaîne reproductible : préparation $\to$ modèle $\to$ prédiction.
\end{description}
\end{kbox}
\end{column}
\begin{column}{0.47\textwidth}
\begin{wbox}[Trois pièges qui coûtent des mois en production]
\begin{enumerate}
\item \textbf{Fuite de données} : donner au modèle une information
du futur (ex. inclure l'alarme de 14 h pour prédire l'incident de 14 h\ldots
qui est \emph{causée} par le pic). On y revient en séance 4.
\item \textbf{Validation mélangée} : découper des séries temporelles au hasard
--- le train contient alors le futur du test.
\item \textbf{Métrique mal choisie} : exactitude sur des incidents rares
(0,2 \% du temps) $\Rightarrow$ un modèle nul paraît parfait.
\end{enumerate}
\end{wbox}
\end{column}
\end{columns}
\end{frame}

\begin{frame}{Quiz 1 --- les fondations (5 min)}
\begin{qbox}[3 questions --- répondez sans regarder la suite]
\begin{enumerate}
\item Détecter des sessions d'authentification anormales \textbf{sans} exemples
d'attaques passés est un problème\ldots\\
(a) supervisé \quad (b) non supervisé \quad (c) par renforcement
\item Un modèle atteint 99,98 \% d'exactitude sur la détection d'incidents
(16 blocs / mois). Que conclure ?\\
(a) excellent modèle \quad (b) il dit peut-être toujours « non » --- il faut
la matrice de confusion \quad (c) il faut plus de CPU
\item Un SLI se mesure, un SLO est\ldots\\
(a) le contrat signé avec le client \quad (b) l'objectif interne qu'on se fixe
sur le SLI \quad (c) une métrique système
\end{enumerate}
\end{qbox}
\vspace{1mm}
\begin{solbox}[Réponses]
\textbf{1 : b} --- sans étiquettes, on cherche l'improbabilité : anomalies (séance 5).
\quad
\textbf{2 : b} --- 99,98 \% = ne jamais alerter ; l'exactitude est piégeuse sur classes rares.
\quad
\textbf{3 : b} --- le SLA est le contrat (avec pénalités), le SLO l'objectif interne.
\end{solbox}
\end{frame}
